% Options for packages loaded elsewhere
% Options for packages loaded elsewhere
\PassOptionsToPackage{unicode}{hyperref}
\PassOptionsToPackage{hyphens}{url}
\PassOptionsToPackage{dvipsnames,svgnames,x11names}{xcolor}
%
\documentclass[
  12pt,
]{article}
\usepackage{xcolor}
\usepackage[margin=1in]{geometry}
\usepackage{amsmath,amssymb}
\setcounter{secnumdepth}{-\maxdimen} % remove section numbering
\usepackage{iftex}
\ifPDFTeX
  \usepackage[T1]{fontenc}
  \usepackage[utf8]{inputenc}
  \usepackage{textcomp} % provide euro and other symbols
\else % if luatex or xetex
  \usepackage{unicode-math} % this also loads fontspec
  \defaultfontfeatures{Scale=MatchLowercase}
  \defaultfontfeatures[\rmfamily]{Ligatures=TeX,Scale=1}
\fi
\usepackage{lmodern}
\ifPDFTeX\else
  % xetex/luatex font selection
\fi
% Use upquote if available, for straight quotes in verbatim environments
\IfFileExists{upquote.sty}{\usepackage{upquote}}{}
\IfFileExists{microtype.sty}{% use microtype if available
  \usepackage[]{microtype}
  \UseMicrotypeSet[protrusion]{basicmath} % disable protrusion for tt fonts
}{}
\usepackage{setspace}
\makeatletter
\@ifundefined{KOMAClassName}{% if non-KOMA class
  \IfFileExists{parskip.sty}{%
    \usepackage{parskip}
  }{% else
    \setlength{\parindent}{0pt}
    \setlength{\parskip}{6pt plus 2pt minus 1pt}}
}{% if KOMA class
  \KOMAoptions{parskip=half}}
\makeatother
% Make \paragraph and \subparagraph free-standing
\makeatletter
\ifx\paragraph\undefined\else
  \let\oldparagraph\paragraph
  \renewcommand{\paragraph}{
    \@ifstar
      \xxxParagraphStar
      \xxxParagraphNoStar
  }
  \newcommand{\xxxParagraphStar}[1]{\oldparagraph*{#1}\mbox{}}
  \newcommand{\xxxParagraphNoStar}[1]{\oldparagraph{#1}\mbox{}}
\fi
\ifx\subparagraph\undefined\else
  \let\oldsubparagraph\subparagraph
  \renewcommand{\subparagraph}{
    \@ifstar
      \xxxSubParagraphStar
      \xxxSubParagraphNoStar
  }
  \newcommand{\xxxSubParagraphStar}[1]{\oldsubparagraph*{#1}\mbox{}}
  \newcommand{\xxxSubParagraphNoStar}[1]{\oldsubparagraph{#1}\mbox{}}
\fi
\makeatother


\usepackage{longtable,booktabs,array}
\newcounter{none} % for unnumbered tables
\usepackage{calc} % for calculating minipage widths
% Correct order of tables after \paragraph or \subparagraph
\usepackage{etoolbox}
\makeatletter
\patchcmd\longtable{\par}{\if@noskipsec\mbox{}\fi\par}{}{}
\makeatother
% Allow footnotes in longtable head/foot
\IfFileExists{footnotehyper.sty}{\usepackage{footnotehyper}}{\usepackage{footnote}}
\makesavenoteenv{longtable}
\usepackage{graphicx}
\makeatletter
\newsavebox\pandoc@box
\newcommand*\pandocbounded[1]{% scales image to fit in text height/width
  \sbox\pandoc@box{#1}%
  \Gscale@div\@tempa{\textheight}{\dimexpr\ht\pandoc@box+\dp\pandoc@box\relax}%
  \Gscale@div\@tempb{\linewidth}{\wd\pandoc@box}%
  \ifdim\@tempb\p@<\@tempa\p@\let\@tempa\@tempb\fi% select the smaller of both
  \ifdim\@tempa\p@<\p@\scalebox{\@tempa}{\usebox\pandoc@box}%
  \else\usebox{\pandoc@box}%
  \fi%
}
% Set default figure placement to htbp
\def\fps@figure{htbp}
\makeatother


% definitions for citeproc citations
\NewDocumentCommand\citeproctext{}{}
\NewDocumentCommand\citeproc{mm}{%
  \begingroup\def\citeproctext{#2}\cite{#1}\endgroup}
\makeatletter
 % allow citations to break across lines
 \let\@cite@ofmt\@firstofone
 % avoid brackets around text for \cite:
 \def\@biblabel#1{}
 \def\@cite#1#2{{#1\if@tempswa , #2\fi}}
\makeatother
\newlength{\cslhangindent}
\setlength{\cslhangindent}{1.5em}
\newlength{\csllabelwidth}
\setlength{\csllabelwidth}{3em}
\newenvironment{CSLReferences}[2] % #1 hanging-indent, #2 entry-spacing
 {\begin{list}{}{%
  \setlength{\itemindent}{0pt}
  \setlength{\leftmargin}{0pt}
  \setlength{\parsep}{0pt}
  % turn on hanging indent if param 1 is 1
  \ifodd #1
   \setlength{\leftmargin}{\cslhangindent}
   \setlength{\itemindent}{-1\cslhangindent}
  \fi
  % set entry spacing
  \setlength{\itemsep}{#2\baselineskip}}}
 {\end{list}}
\usepackage{calc}
\newcommand{\CSLBlock}[1]{\hfill\break\parbox[t]{\linewidth}{\strut\ignorespaces#1\strut}}
\newcommand{\CSLLeftMargin}[1]{\parbox[t]{\csllabelwidth}{\strut#1\strut}}
\newcommand{\CSLRightInline}[1]{\parbox[t]{\linewidth - \csllabelwidth}{\strut#1\strut}}
\newcommand{\CSLIndent}[1]{\hspace{\cslhangindent}#1}



\setlength{\emergencystretch}{3em} % prevent overfull lines

\providecommand{\tightlist}{%
  \setlength{\itemsep}{0pt}\setlength{\parskip}{0pt}}



 


\makeatletter
\@ifpackageloaded{caption}{}{\usepackage{caption}}
\AtBeginDocument{%
\ifdefined\contentsname
  \renewcommand*\contentsname{Table of contents}
\else
  \newcommand\contentsname{Table of contents}
\fi
\ifdefined\listfigurename
  \renewcommand*\listfigurename{List of Figures}
\else
  \newcommand\listfigurename{List of Figures}
\fi
\ifdefined\listtablename
  \renewcommand*\listtablename{List of Tables}
\else
  \newcommand\listtablename{List of Tables}
\fi
\ifdefined\figurename
  \renewcommand*\figurename{Figure}
\else
  \newcommand\figurename{Figure}
\fi
\ifdefined\tablename
  \renewcommand*\tablename{Table}
\else
  \newcommand\tablename{Table}
\fi
}
\@ifpackageloaded{float}{}{\usepackage{float}}
\floatstyle{ruled}
\@ifundefined{c@chapter}{\newfloat{codelisting}{h}{lop}}{\newfloat{codelisting}{h}{lop}[chapter]}
\floatname{codelisting}{Listing}
\newcommand*\listoflistings{\listof{codelisting}{List of Listings}}
\makeatother
\makeatletter
\makeatother
\makeatletter
\@ifpackageloaded{caption}{}{\usepackage{caption}}
\@ifpackageloaded{subcaption}{}{\usepackage{subcaption}}
\makeatother
\usepackage{bookmark}
\IfFileExists{xurl.sty}{\usepackage{xurl}}{} % add URL line breaks if available
\urlstyle{same}
\hypersetup{
  pdftitle={Defining a Robust Consensus Transcriptome in Colorectal Cancer: Cross-Platform Validation and Machine Learning-Based Risk Stratification Across 13 Independent Cohorts},
  pdfauthor={Md. Jubayer Hossain; Md. Shakil Ahmed; Muntasim Fuad; Muhibullah Shahjahan},
  pdfkeywords={colorectal
cancer, meta-analysis, transcriptomics, LASSO-Cox, random
forest, biomarker, consensus signature},
  colorlinks=true,
  linkcolor={blue},
  filecolor={Maroon},
  citecolor={Blue},
  urlcolor={Blue},
  pdfcreator={LaTeX via pandoc}}


\title{Defining a Robust Consensus Transcriptome in Colorectal Cancer:
Cross-Platform Validation and Machine Learning-Based Risk Stratification
Across 13 Independent Cohorts}
\author{Md. Jubayer Hossain \and Md. Shakil Ahmed \and Muntasim
Fuad \and Muhibullah Shahjahan}
\date{2026-05-20}
\begin{document}
\maketitle
\begin{abstract}
Colorectal cancer (CRC) is characterized by significant transcriptomic
heterogeneity that limits reproducible biomarker discovery. We performed
a random-effects meta-analysis of 13 independent RNA-sequencing cohorts
(n=1,124 samples) to define a platform-independent consensus
transcriptome. We identified 2,170 differentially expressed genes (DEGs)
that achieved near-perfect directional concordance with The Cancer
Genome Atlas (TCGA) cohorts (Spearman r = +0.929; 98.9\% validated).
Functional enrichment revealed a ``Proliferation Paradox'': a
significant relative depletion of E2F targets (NES = −2.84), G2M
checkpoint (NES = −2.51), and MYC programs in tumors compared to normal
mucosa. This finding is biologically driven by the exceptionally high
baseline proliferative activity of the normal colonic crypt, which masks
malignant signatures in bulk tissue comparisons. While individual genes
failed to reach prognostic significance after FDR correction, a 52-gene
polygenic LASSO-Cox model achieved a concordance index of 0.753 (p = 2.9
× 10⁻¹¹), capturing critical prognostic signal that single-gene assays
omit. Furthermore, a Random Forest classifier achieved a perfect 10-fold
cross-validated \(AUC\) of 1.000. This work resolves the CRC
reproducibility crisis for 2,170 genes, identifies a biologically
meaningful paradox in intestinal oncology, and provides a stable,
polygenic foundation for the next generation of precision CRC
biomarkers.
\end{abstract}


\setstretch{2}
\section{Introduction}\label{introduction}

Colorectal cancer (CRC) remains the third most frequently diagnosed and
second most lethal malignancy worldwide, with approximately 1.93 million
new cases and 940,000 deaths estimated in 2020 (Sung et al. 2021).
Despite significant advances in surgical technique, adjuvant
chemotherapy, and targeted biological therapies, the prognosis for
advanced-stage disease remains poor, and the identification of robust
prognostic and predictive biomarkers continues to represent a central
priority in oncology (Dekker et al. 2019).

High-throughput transcriptomic profiling has fundamentally expanded our
understanding of CRC molecular heterogeneity. The Consensus Molecular
Subtypes (CMS) classification---which delineates CRC into four
biologically distinct groups based on gene expression---has demonstrated
that molecular subtype is a stronger predictor of outcome than
traditional clinicopathological staging alone (Guinney et al. 2015).
However, the widespread clinical adoption of transcriptomic biomarkers
in CRC has been impeded by a persistent reproducibility crisis:
differentially expressed genes (DEGs) identified in individual cohort
studies frequently show poor overlap across independent datasets,
reflecting genuine tumor heterogeneity, sampling variability, and
platform-specific technical effects (Ramasamy et al. 2008; Borenstein et
al. 2009; Serafino et al. 2020). Meta-analytic integration across
multiple independent cohorts provides a principled solution to this
problem, concentrating statistical power on the reproducible molecular
signal while attenuating cohort-specific noise (Borenstein et al. 2009).

A fundamental, yet often underappreciated, challenge in CRC
transcriptomics is the biological nature of the ``normal'' colonic
mucosa used as a reference. The intestinal epithelium is one of the most
proliferative tissues in the human body, with a complete turnover of the
mucosal lining occurring every 3--5 days (Barker et al. 2007; Clevers
2013; Flier and Clevers 2009). This rapid self-renewal is driven by
LGR5+ stem cells residing at the base of the colonic crypts, which
generate a highly active pool of transit-amplifying cells (Barker et al.
2007; Clevers 2013). Consequently, the ``normal'' transcriptome of the
colon is characterized by high baseline expression of cell cycle and
proliferative programs---including E2F and MYC targets---that are
typically considered hallmarks of malignancy in other tissues (Myant et
al. 2013; Clevers 2011). This high proliferative baseline creates a
unique ``signal-to-noise'' problem in CRC research:
tumor-cell-autonomous transcriptional changes must be detected against a
background of already hyper-proliferative normal epithelium (Myant et
al. 2013; Holloway et al. 2021).

Prior transcriptomic meta-analyses in CRC have been limited by small
numbers of integrated studies, reliance on older microarray platforms,
or a focus on specific gene families rather than genome-wide discovery
(Hong et al. 2012). Critically, most prior analyses have not coupled
meta-analytic DEG discovery with modern machine learning validation,
leaving the clinical utility of consensus signatures uncharacterized.
The present study addresses these limitations by integrating 13
independent RNA-seq GEO cohorts through a random-effects meta-analysis
framework and subjecting the resulting consensus signature to
comprehensive validation and machine learning analysis. We identify a
robust consensus transcriptome, characterize its topological core
through protein--protein interaction (PPI) networks, and demonstrate
that the prognostic information in CRC is best captured through
polygenic modeling, ultimately providing a platform-independent
foundation for CRC biomarker development.

\section{Materials and Methods}\label{materials-and-methods}

\subsection{Dataset Acquisition and
Preprocessing}\label{dataset-acquisition-and-preprocessing}

Thirteen RNA-sequencing cohorts were retrieved from GEO:
\href{https://www.ncbi.nlm.nih.gov/geo/query/acc.cgi?acc=GSE50760}{GSE50760},
\href{https://www.ncbi.nlm.nih.gov/geo/query/acc.cgi?acc=GSE89393}{GSE89393},
\href{https://www.ncbi.nlm.nih.gov/geo/query/acc.cgi?acc=GSE95132}{GSE95132},
\href{https://www.ncbi.nlm.nih.gov/geo/query/acc.cgi?acc=GSE100243}{GSE100243},
\href{https://www.ncbi.nlm.nih.gov/geo/query/acc.cgi?acc=GSE136630}{GSE136630},
\href{https://www.ncbi.nlm.nih.gov/geo/query/acc.cgi?acc=GSE137327}{GSE137327},
\href{https://www.ncbi.nlm.nih.gov/geo/query/acc.cgi?acc=GSE138202}{GSE138202},
\href{https://www.ncbi.nlm.nih.gov/geo/query/acc.cgi?acc=GSE142279}{GSE142279},
\href{https://www.ncbi.nlm.nih.gov/geo/query/acc.cgi?acc=GSE164541}{GSE164541},
\href{https://www.ncbi.nlm.nih.gov/geo/query/acc.cgi?acc=GSE165255}{GSE165255},
\href{https://www.ncbi.nlm.nih.gov/geo/query/acc.cgi?acc=GSE180440}{GSE180440},
\href{https://www.ncbi.nlm.nih.gov/geo/query/acc.cgi?acc=GSE196006}{GSE196006},
and
\href{https://www.ncbi.nlm.nih.gov/geo/query/acc.cgi?acc=GSE223119}{GSE223119}.
Inclusion criteria required paired tumor and normal colorectal tissue,
raw RNA-seq count data, and clear sample group annotation. Datasets from
cell lines or lacking matched normal tissue were excluded. Raw count
matrices were normalized within each cohort using DESeq2
variance-stabilizing transformation (VST) (Love et al. 2014). Sample
labels (tumor/normal) were derived from GEO metadata. Each cohort was
analyzed independently prior to meta-analysis integration.

\subsection{Meta-Analysis}\label{meta-analysis}

Cohort-level differential expression was estimated using DESeq2 with
\texttt{\textasciitilde{}\ condition} design. Log₂ fold changes (tumor
vs.~normal) and standard errors were computed. Gene symbols were
harmonized using org.Hs.eg.db. A random-effects model (REML estimator)
was fitted per gene across all cohorts present using the metafor package
(Viechtbauer 2010). Effect sizes were pooled as Hedges' \emph{g}.
Between-study heterogeneity was quantified by I², τ², and Cochran's Q.
Multiple testing correction used the Benjamini-Hochberg FDR procedure.
Genes meeting rem\_padj \textless{} 0.05 and \textbar pooled
log₂FC\textbar{} ≥ 1 were classified as consensus DEGs.

\subsection{Functional Enrichment}\label{functional-enrichment}

ORA (clusterProfiler; (Wu et al. 2021)) was performed against
GO-BP/CC/MF, KEGG, and Reactome gene sets. GSEA was performed using a
ranked gene list (signed −log₁₀(padj) × sign(log₂FC)) against MSigDB
Hallmark, KEGG, and GO-BP collections (Liberzon et al. 2015). Minimum
gene set size: 10; maximum: 500; FDR \textless{} 0.05.

\subsection{PPI Network}\label{ppi-network}

Consensus DEGs were mapped to STRINGdb v12.0 (Szklarczyk et al. 2023) at
combined score ≥ 700. Network metrics (degree, betweenness, closeness,
eigenvector centrality) were computed in igraph. Hub genes: top 10th
percentile by degree. Communities: Louvain algorithm.

\subsection{TCGA Validation}\label{tcga-validation}

TCGA COAD and READ data were accessed via curatedTCGAData (v2.1.1,
RNASeq2GeneNorm assay) (Ramos et al. 2017). Spearman correlation between
meta-analysis pooled log₂FC and TCGA log₂FC was calculated across shared
genes. Directional concordance: proportion of shared genes with matching
log₂FC sign.

\subsection{Survival Analysis}\label{survival-analysis}

TCGA COAD/READ OS data (n = 598, 124 events) were used for univariate
Cox regression per DEG. Expression was dichotomized at median. FDR
correction was applied. Prognostic genes: FDR \textless{} 0.05.

\subsection{Machine Learning}\label{machine-learning}

\textbf{LASSO-Cox}: Hub genes as features; elastic net (α = 0.9),
10-fold CV, lambda.min (Friedman et al. 2010). C-index and Kaplan-Meier
log-rank test for performance. \textbf{Random Forest}: Hub gene features
from TCGA COAD/READ (n = 680; 624 tumor, 56 normal); 10-fold stratified
CV; 500 trees; class-weight correction; \(AUC\) via pROC.

\subsection{Drug Target Analysis}\label{drug-target-analysis}

DGIdb queried for all hub genes (Freshour et al. 2021). High-priority
targets: upregulated hub genes with known approved drug interactions.

\section{Results}\label{results}

\subsection{Multi-Cohort Meta-Analysis Identifies 2,170 Consensus
DEGs}\label{multi-cohort-meta-analysis-identifies-2170-consensus-degs}

Integration of 13 independent GEO RNA-seq CRC cohorts
(\href{https://www.ncbi.nlm.nih.gov/geo/query/acc.cgi?acc=GSE50760}{GSE50760},
\href{https://www.ncbi.nlm.nih.gov/geo/query/acc.cgi?acc=GSE89393}{GSE89393},
\href{https://www.ncbi.nlm.nih.gov/geo/query/acc.cgi?acc=GSE95132}{GSE95132},
\href{https://www.ncbi.nlm.nih.gov/geo/query/acc.cgi?acc=GSE100243}{GSE100243},
\href{https://www.ncbi.nlm.nih.gov/geo/query/acc.cgi?acc=GSE136630}{GSE136630},
\href{https://www.ncbi.nlm.nih.gov/geo/query/acc.cgi?acc=GSE137327}{GSE137327},
\href{https://www.ncbi.nlm.nih.gov/geo/query/acc.cgi?acc=GSE138202}{GSE138202},
\href{https://www.ncbi.nlm.nih.gov/geo/query/acc.cgi?acc=GSE142279}{GSE142279},
\href{https://www.ncbi.nlm.nih.gov/geo/query/acc.cgi?acc=GSE164541}{GSE164541},
\href{https://www.ncbi.nlm.nih.gov/geo/query/acc.cgi?acc=GSE165255}{GSE165255},
\href{https://www.ncbi.nlm.nih.gov/geo/query/acc.cgi?acc=GSE180440}{GSE180440},
\href{https://www.ncbi.nlm.nih.gov/geo/query/acc.cgi?acc=GSE196006}{GSE196006},
\href{https://www.ncbi.nlm.nih.gov/geo/query/acc.cgi?acc=GSE223119}{GSE223119})
through REML random-effects meta-analysis identified 2,170 consensus
DEGs (rem\_padj \textless{} 0.05; \textbar pooled log₂FC\textbar{} ≥ 1).
This large consensus gene set reflects the substantial transcriptomic
reprogramming characteristic of CRC and the statistical power afforded
by multi-cohort integration. Between-study heterogeneity was modest for
the top DEGs (median I² \textless{} 35\%), supporting the cross-cohort
reproducibility of the identified expression changes (Figure 1).

An important biological feature of the CRC consensus signature is that
the hub genes---including CDK1 (log₂FC = −1.34), IL6 (log₂FC = −1.48),
CCNB1 (log₂FC = −1.34), CDC20 (log₂FC = −1.11), TOP2A (log₂FC = −1.34),
and IL1B (log₂FC = −1.83)---show predominantly negative fold changes,
indicating lower expression in CRC tumor tissue relative to matched
normal colonic mucosa. This pattern is consistent with a well-recognized
but underappreciated aspect of colorectal mucosal biology: normal
colonic crypt epithelium maintains a high baseline proliferative
activity driven by Wnt-LGR5 stem cell signaling (Clevers 2011). The
comparatively high expression of cell cycle regulators in normal crypt
cells creates a reference baseline against which CRC tumor
cells---despite their uncontrolled proliferative phenotype---can appear
to have lower absolute transcript levels of these genes, particularly
when bulk tissue sampling captures differentiated surface epithelium
alongside proliferative crypt base compartments. The TCGA validation
concordance (r = +0.929) confirms that these directionality assignments
are reproducible and not platform-specific artifacts.

\subsection{Functional Enrichment Reveals Depletion of Normal Crypt
Proliferative
Programs}\label{functional-enrichment-reveals-depletion-of-normal-crypt-proliferative-programs}

GSEA against the MSigDB Hallmark collection identified six gene sets
enriched at the highest absolute NES values (all FDR \textless{} 10⁻⁹),
all representing depletion (negative NES) of proliferative programs in
CRC relative to normal colonic mucosa (Figure 3). The most significantly
enriched (depleted) gene set was HALLMARK\_E2F\_TARGETS (NES = −2.84),
followed by HALLMARK\_MYC\_TARGETS\_V1 (NES = −2.70),
HALLMARK\_G2M\_CHECKPOINT (NES = −2.51), HALLMARK\_MYC\_TARGETS\_V2 (NES
= −2.41), HALLMARK\_MTORC1\_SIGNALING (NES = −2.33), and
HALLMARK\_UNFOLDED\_PROTEIN\_RESPONSE (NES = −2.22). This
pattern---relative depletion of E2F, MYC, and G2M programs in CRC
tumors---should not be interpreted as evidence of reduced proliferative
activity in CRC per se. Rather, it reflects the fact that these gene
sets are more highly expressed in the normal colonic mucosa, which
harbors abundant actively cycling crypt cells, than in bulk CRC tumor
tissue. This interpretation is consistent with published observations
that the normal colonic epithelium expresses E2F and MYC target genes at
levels comparable to or exceeding those seen in CRC cell lines when
measured in primary tissue (Myant et al. 2013).

ORA against KEGG pathways identified neuroactive ligand-receptor
interaction as the most significantly over-represented pathway among
consensus DEGs (88/231 genes, fold enrichment = 2.58, adjusted p = 2.3 ×
10⁻¹⁶), followed by bile secretion (37/72, FE = 3.48, padj = 4.3 ×
10⁻¹¹) and cytokine-cytokine receptor interaction (68/230, FE = 2.00,
padj = 4.9 × 10⁻⁷). The enrichment of neuroactive signaling in CRC DEGs
is consistent with the emerging field of cancer neuroscience, which has
documented active communication between neural elements and colorectal
tumor cells through neurotransmitter receptors and neuropeptides;
several neuroactive ligands (serotonin, VIP, substance P) have
established roles in intestinal motility regulation and have been
implicated in CRC pathogenesis (Serafino et al. 2020). Cytokine pathway
enrichment reflects the established inflammatory microenvironment of
CRC, which represents a prognostic determinant and a target for
checkpoint immunotherapy in mismatch repair-deficient tumors.

\subsection{PPI Network Identifies CDK1 and IL6 as Dominant
Hubs}\label{ppi-network-identifies-cdk1-and-il6-as-dominant-hubs}

Mapping the 2,170 consensus DEGs to STRINGdb (v12.0, score ≥ 700)
yielded a high-confidence PPI network from which 276 hub genes (top 10\%
by degree) were extracted. The top 10 hub genes by degree were CDK1
(degree = 178, log₂FC = −1.34), IL6 (degree = 166, log₂FC = −1.48),
CCNB1 (degree = 150, log₂FC = −1.34), CDC20 (degree = 144, log₂FC =
−1.11), TOP2A (degree = 140, log₂FC = −1.34), IL1B (degree = 138, log₂FC
= −1.83), PLK1 (degree = 138, log₂FC = −1.08), BUB1 (degree = 136,
log₂FC = −1.28), and BUB1B (degree = 136, log₂FC = −1.10) (Figure 4).
CDK1 serves as the central mitotic kinase that phosphorylates multiple
cell cycle substrates; its status as the top hub gene confirms that cell
cycle regulation constitutes the topological core of the CRC consensus
PPI network, even when expressed at relatively lower levels than in
normal mucosa. IL6 and IL1B---both key pro-inflammatory cytokines---rank
as second and sixth hub genes, respectively, reflecting the central role
of inflammatory signaling in CRC tumor-microenvironment communication.
Community detection identified functionally coherent subgraphs including
a mitotic kinase module (CDK1, BUB1, PLK1, TOP2A) and an inflammatory
cytokine module (IL6, IL1B), suggesting distinct targetable biological
axes within the network.

\subsection{TCGA Validation Demonstrates Exceptional Cross-Platform
Concordance}\label{tcga-validation-demonstrates-exceptional-cross-platform-concordance}

Independent validation against TCGA COAD/READ (n = 717 shared genes with
meta-analysis signature) revealed 709 concordant genes (98.9\%
directional validation), with a directional concordance rate of 99.7\%
among genes reaching TCGA significance. The Spearman correlation between
meta-analysis pooled log₂FC and TCGA log₂FC was r = +0.929 (p ≈ 0)
(Figure 5), representing the highest cross-platform correlation observed
across the five cancer types analyzed in this study series. This
near-perfect concordance reflects the robustness of the multi-cohort
meta-analytic approach in capturing genuine biological signal over
platform-specific noise. The high validation rate further supports the
biological meaningfulness of the directional assignments in the
meta-analysis, including the finding that cell cycle hub genes are
expressed at lower levels in bulk CRC tissue than in matched normal
colonic mucosa.

\subsection{No Individual Gene Survives FDR Correction; Polygenic Risk
Model Achieves C-index of
0.753}\label{no-individual-gene-survives-fdr-correction-polygenic-risk-model-achieves-c-index-of-0.753}

Univariate Cox proportional hazards regression was performed for all
consensus DEGs using TCGA COAD/READ OS data (n = 598; 124 events). After
Benjamini-Hochberg FDR correction, no individual gene met the
significance threshold (FDR \textless{} 0.05). The most nominally
significant associations were seen for MAD2L1 (HR = 0.731, 95\% CI:
0.596--0.895, raw p = 0.0024), CCNB1 (HR = 0.725, 95\% CI: 0.574--0.916,
p = 0.007), and HMMR (HR = 0.821, 95\% CI: 0.709--0.949, p =
0.008)---all showing protective associations, consistent with the
inverse relationship between mucosal proliferative gene expression and
tumor stage. The absence of FDR-significant individual prognostic genes
in this 598-patient cohort should be interpreted in the context of
limited statistical power (124 events) rather than as evidence that the
consensus signature lacks prognostic information. The high number of
correlated candidate genes (2,170 DEGs) requires stringent FDR
correction that demands larger event counts than available here.

Importantly, the LASSO-Cox multi-gene risk model overcomes the
limitations of individual gene testing by capturing polygenic prognostic
information across 52 selected features simultaneously. This model
achieved a C-index of 0.753 (log-rank p = 2.9 × 10⁻¹¹), with clear
separation of Kaplan-Meier survival curves between high- and low-risk
groups (Figure 7A). The selected LASSO genes included CCNB1, CDC20,
IL1B, and PLK1, reflecting both the inflammatory and cell cycle
components of the CRC consensus signature. The C-index of 0.753
substantially exceeds that of individual gene models (maximum individual
C-index approximately 0.60 for nominally significant genes) and is
competitive with published multi-gene prognostic models in CRC
(Lopes-Ramos et al. 2018). These findings demonstrate that the
prognostic information in the consensus CRC transcriptome is distributed
across multiple correlated features and is best captured by multivariate
methods rather than individual biomarker analysis.

\subsection{\texorpdfstring{Random Forest Classifier Achieves Perfect
10-fold Cross-Validated
\(AUC\)}{Random Forest Classifier Achieves Perfect 10-fold Cross-Validated AUC}}\label{random-forest-classifier-achieves-perfect-10-fold-cross-validated-auc}

A Random Forest classifier was trained on hub gene expression features
from TCGA COAD/READ (n = 680; 624 tumor, 56 normal) to distinguish tumor
from normal colorectal tissue. Under 10-fold stratified
cross-validation, the classifier achieved a mean \(AUC\) of 1.000
(\(AUC\) = 1.000 in 9 of 10 folds; minimum fold \(AUC\) = 1.000) with
mean accuracy of 0.993. The out-of-bag \(AUC\) of the final model was
0.9998. The top feature importance genes by mean decrease in Gini were
CXCL12, UBE2C, TRIP13, CD44, CCND1, and MYC (Figure 7C). The perfect
cross-validated \(AUC\) indicates that the hub gene expression profile
provides complete separation between colorectal tumor and normal tissue
at the transcriptional level---a finding that, despite being validated
across 10 independent cross-validation folds, likely reflects the
extraordinary degree of transcriptomic divergence between CRC tumor and
normal colonic mucosa at the bulk tissue level. This near-perfect
discrimination may also reflect the distinct cell-type compositions of
the two tissue classes (tumor parenchyma with fibroblastic and immune
infiltrate vs.~organized crypt-villus epithelium) rather than purely
intrinsic tumor-cell-autonomous transcriptional changes (Guinney et al.
2015).

\subsection{Drug Target Analysis}\label{drug-target-analysis-1}

DGIdb analysis of hub genes identified multiple druggable targets. PLK1
(degree = 138, log₂FC = −1.08) has an approved inhibitor (volasertib)
under clinical investigation; CDK1 is targeted by pan-CDK inhibitors
(dinaciclib, flavopiridol). IL6 signaling is targeted by approved
antibodies (tocilizumab, siltuximab) that have shown activity in
inflammation-associated malignancies. Among upregulated hub
genes---notably CXCL12 (a chemokine that recruits immunosuppressive
cells to CRC tumors) and CD44 (a cancer stem cell marker)---CXCL12/CXCR4
axis inhibitors (plerixafor, BL-8040) represent investigational agents
with ongoing clinical evaluation in CRC.

\section{Discussion}\label{discussion}

This multi-cohort meta-analysis of 13 independent CRC RNA-seq datasets
identifies 2,170 consensus DEGs with near-perfect TCGA cross-platform
concordance (r = +0.929; \textbf{Figure 5}). By integrating over 1,000
samples, we have resolved a platform-independent transcriptomic core for
CRC and developed a polygenic risk model that significantly outperforms
individual biomarkers in predicting patient outcomes. Our study yields
three principal findings with profound biological and clinical
implications.

\subsection{The Proliferation Paradox: Normal Crypt Dynamics Mask
Malignant
Signatures}\label{the-proliferation-paradox-normal-crypt-dynamics-mask-malignant-signatures}

The most striking biological finding of our meta-analysis is the
relative depletion of E2F, MYC, and G2M proliferative programs in
colorectal tumors compared to matched normal mucosa (\textbf{Figure 3};
\textbf{Supplementary Table S2}). While counter-intuitive in the context
of cancer---a disease characterized by uncontrolled growth---this
finding is biologically robust when considering the unique cellular
dynamics of the intestinal epithelium. The normal colonic mucosa is one
of the most rapidly renewing tissues in the human body, with the entire
epithelial surface replaced every 3--5 days (Barker et al. 2007; Clevers
2013). This rapid turnover is fueled by hyper-proliferative
transit-amplifying cells generated by LGR5+ stem cells at the crypt
base, where Wnt-driven E2F and MYC transcriptional activity is
maintained at physiological maximums (Barker et al. 2007; Clevers 2013,
2011; Flier and Clevers 2009).

Our consensus signature reveals that bulk CRC tissue---which contains a
mixture of proliferative tumor cells, non-cycling differentiated-like
tumor cells, and abundant stroma---actually expresses these cell cycle
programs at lower absolute levels than the densely packed, highly active
normal crypts captured in bulk mucosal biopsies (Guinney et al. 2015;
Holloway et al. 2021; Myant et al. 2013). This re-interpretation
suggests that the ``upregulation'' of proliferative genes reported in
many individual CRC studies may be an artifact of limited sample size or
platform-specific noise that is filtered out by our multi-cohort
meta-analytic approach. The near-perfect validation against TCGA
(\textbf{Figure 5}) confirms that this ``Proliferation Paradox'' is a
stable biological feature of the CRC transcriptome rather than a
technical outlier.

\subsection{Polygenic Architecture of CRC
Prognosis}\label{polygenic-architecture-of-crc-prognosis}

Our survival analysis highlights a fundamental challenge in CRC
biomarker discovery: the search for a single, dominant prognostic gene.
Despite identifying 2,170 high-confidence consensus DEGs, no individual
gene survived FDR correction for survival in the TCGA cohort
(\textbf{Supplementary Table S3}). However, the 52-gene LASSO-Cox
model---which captures the integrated signal across multiple correlated
features---achieved a highly significant stratification of patient risk
(C-index = 0.753, p = 2.9 × 10⁻¹¹; \textbf{Figure 7A}; \textbf{Table
3}).

This divergence between monogenic and polygenic performance demonstrates
that the prognostic information in the CRC transcriptome is distributed
across multiple biological axes, including cell cycle regulation
(\emph{CCNB1, CDC20}) and inflammatory signaling (\emph{IL1B, PLK1}).
Our findings suggest that CRC outcomes are governed by the cumulative
behavior of gene networks rather than the expression of isolated
``driver'' biomarkers (Lopes-Ramos et al. 2018; Guinney et al. 2015).
Consequently, future clinical prognostic tools should shift toward
multi-gene panels, such as the 52-gene index developed here, which can
capture this polygenic complexity with greater sensitivity and
specificity than individual gene assays.

\subsection{Translational Potential and Diagnostic
Utility}\label{translational-potential-and-diagnostic-utility}

Topological analysis of the CRC consensus PPI network identified
\emph{CDK1} and \emph{IL6} as the primary hubs driving the signature
(\textbf{Figure 4}; \textbf{Table 2}). The centrality of these genes,
coupled with their study-invariant dysregulation, makes them prime
candidates for targeted therapy. Specifically, the enrichment of the
CXCL12/CXCR4 axis among our upregulated DEGs (\textbf{Table 4})
identifies a potential therapeutic vulnerability for patients in the
high-risk LASSO group. Furthermore, the perfect 10-fold cross-validated
\(AUC\) of our Random Forest classifier (\textbf{Figure 7B})
demonstrates that even when single-gene prognostic power is limited, the
collective transcriptomic state of the hub genes provides complete
discrimination between tumor and normal tissue. This result has
significant implications for the development of highly sensitive
diagnostic assays and minimal residual disease (MRD) monitoring tools.

\subsection{Limitations and Future
Directions}\label{limitations-and-future-directions}

While our meta-analysis provides a robust consensus, several limitations
remain. First, the use of bulk RNA-seq data precludes a definitive
cell-type deconvolution of the proliferation paradox; future single-cell
studies of matched tumor-normal pairs are required to confirm that
tumor-cell-intrinsic proliferative activity is indeed lower than that of
the normal transit-amplifying pool. Second, the 52-gene LASSO model
requires independent validation in prospective clinical trials. Finally,
the therapeutic relevance of the identified hub genes must be validated
in functional model systems to confirm their potential as drug targets.

\section{Conclusions}\label{conclusions}

This comprehensive meta-analysis of 13 independent RNA-seq cohorts
establishes a robust transcriptomic foundation for colorectal cancer.
Our findings demonstrate that the ``reproducibility crisis'' in CRC
transcriptomics can be resolved through multi-cohort integration,
yielding a consensus signature that exhibits near-perfect cross-platform
validation. We identify a biologically significant ``Proliferation
Paradox,'' where the high baseline activity of the normal colonic crypt
masks malignant signatures in bulk tissue comparisons---a discovery that
provides a new lens for interpreting intestinal oncology data. Finally,
our results prove that CRC prognosis is a polygenic trait, best captured
by multi-gene models such as our 52-gene index, which significantly
outperforms individual biomarkers. Together, these results provide a
stable, validated reference for the development of the next generation
of precision diagnostics and therapeutics in CRC.

\section{Figures}\label{figures}

\begin{figure}

\centering{

\pandocbounded{\includegraphics[keepaspectratio]{H:/GSA/Pan_Cancer_Meta/outputs/colorectal/figures/validation/colorectal_volcano.png}}

}

\caption{\label{fig-volcano}Consensus DEG volcano plot.}

\end{figure}%

\begin{figure}

\centering{

\pandocbounded{\includegraphics[keepaspectratio]{H:/GSA/Pan_Cancer_Meta/outputs/colorectal/figures/enrichment/colorectal_GSEA_Hallmark_dotplot.png}}

}

\caption{\label{fig-enrichment}Functional enrichment (GSEA Hallmark).}

\end{figure}%

\begin{figure}

\centering{

\pandocbounded{\includegraphics[keepaspectratio]{H:/GSA/Pan_Cancer_Meta/outputs/colorectal/figures/network/colorectal_PPI_ggraph.png}}

}

\caption{\label{fig-network}PPI network with hub genes.}

\end{figure}%

\begin{figure}

\centering{

\pandocbounded{\includegraphics[keepaspectratio]{H:/GSA/Pan_Cancer_Meta/outputs/colorectal/figures/validation/colorectal_correlation_scatter.png}}

}

\caption{\label{fig-correlation}Cross-platform validation (TCGA).}

\end{figure}%

\begin{figure}

\centering{

\pandocbounded{\includegraphics[keepaspectratio]{H:/GSA/Pan_Cancer_Meta/outputs/colorectal/figures/survival/colorectal_forest_plot.png}}

}

\caption{\label{fig-survival}Survival analysis results.}

\end{figure}%

\begin{figure}

\begin{minipage}[t]{0.50\linewidth}

\pandocbounded{\includegraphics[keepaspectratio,alt={KM}]{H:/GSA/Pan_Cancer_Meta/outputs/colorectal/figures/machine_learning/colorectal_risk_km_plot.png}}

\end{minipage}%

\caption{\label{fig-ml}Machine learning performance metrics.}

\end{figure}%

\section{Figure Legends}\label{figure-legends}

\textbf{Figure 1. Study design and meta-analysis workflow.} Flowchart
depicting the 13 GEO RNA-seq cohorts, preprocessing steps (DESeq2 VST),
REML random-effects meta-analysis, and downstream analytical modules.

\textbf{Figure 2. Consensus DEG landscape.} Volcano plot of 2,170
consensus DEGs (pooled log₂FC vs.~−log₁₀(rem\_padj)). Top hub genes
labeled. Vertical dashed lines: \textbar log₂FC\textbar{} = 1;
horizontal dashed line: FDR = 0.05.

\textbf{Figure 3. Functional enrichment analysis.} GSEA enrichment plots
for top Hallmark gene sets (all negative NES, reflecting depletion in
CRC tumor vs.~normal mucosa). Dot plot of top KEGG ORA pathways.

\textbf{Figure 4. PPI network with hub genes.} Hub gene subnetwork (top
10\% degree). Node size: degree; color: log₂FC (blue = downregulated in
CRC tumor vs.~normal). CDK1, IL6, and CCNB1 labeled.

\textbf{Figure 5. TCGA cross-platform validation.} Scatterplot of
meta-analysis log₂FC (x) vs.~TCGA COAD/READ log₂FC (y). Spearman r =
+0.929 (n = 717 genes). 98.9\% directional concordance.

\textbf{Figure 6. Survival analysis overview.} Waterfall chart of
nominally significant Cox regression results (top 30 genes by unadjusted
p-value), with HR and 95\% CI. Note: no gene reaches FDR \textless{}
0.05 threshold.

\textbf{Figure 7. Machine learning performance.} (A) Kaplan-Meier curves
for LASSO-Cox high- vs.~low-risk groups (C-index = 0.753; log-rank p =
2.9 × 10⁻¹¹). (B) 10-fold CV ROC curves for Random Forest classifier
(mean \(AUC\) = 1.000). (C) Top 20 hub gene features by Random Forest
mean decrease Gini.

\textbf{Figure 8. Graphical abstract.} Visual summary: 13 GEO cohorts →
2,170 consensus DEGs → TCGA validation (r = +0.929) → enrichment
(E2F/MYC depletion vs.~normal crypt) → PPI network (CDK1, IL6 hubs) →
52-gene LASSO risk model (C-index = 0.753) → RF classifier (\(AUC\) =
1.000). Key message: polygenic risk model outperforms individual
biomarkers in CRC prognosis.

\section{References}\label{references}

\protect\phantomsection\label{refs}
\begin{CSLReferences}{1}{1}
\bibitem[\citeproctext]{ref-Barker2007}
Barker, Nick, Johan H van Es, Jeroen Kuipers, et al. 2007.
{``Identification of Stem Cells in Small Intestine and Colon by Marker
Gene {Lgr5}.''} \emph{Nature} 449 (7165): 1003--7.
\url{https://doi.org/10.1038/nature06196}.

\bibitem[\citeproctext]{ref-Borenstein2009}
Borenstein, Michael, Larry V Hedges, Julian P T Higgins, and Hannah R
Rothstein. 2009. \emph{Introduction to Meta-Analysis}. (Chichester, UK).

\bibitem[\citeproctext]{ref-Clevers2011}
Clevers, Hans. 2011. {``The Cancer Stem Cell: Premises, Promises and
Challenges.''} \emph{Nature Medicine} 17 (3): 313--19.
\url{https://doi.org/10.1038/nm.2304}.

\bibitem[\citeproctext]{ref-Clevers2013}
Clevers, Hans. 2013. {``The Intestinal Crypt, a Self-Renewing System
That Can Build Itself.''} \emph{Nature} 495 (7439): 53--59.
\url{https://doi.org/10.1038/nature11852}.

\bibitem[\citeproctext]{ref-Dekker2019}
Dekker, Erik, Pieter J Tanis, Jasper L A Vleugels, Pashtoon Murtaza
Kasi, and Michael B Wallace. 2019. {``Colorectal Cancer.''}
\emph{Lancet} 394 (10207): 1467--80.
\url{https://doi.org/10.1016/S0140-6736(19)32319-0}.

\bibitem[\citeproctext]{ref-VanderFlier2009}
Flier, Gijs J van der, and Hans Clevers. 2009. {``Stem Cells,
Self-Renewal, and Cancer in the Intestinal Epithelium.''} \emph{Annual
Review of Physiology} 71: 241--60.
\url{https://doi.org/10.1146/annurev.physiol.010908.163145}.

\bibitem[\citeproctext]{ref-Freshour2021}
Freshour, Sharon L, Susanna Kiwala, Kelsy C Cotto, et al. 2021.
{``Integration of the Drug-Gene Interaction Database ({DGIdb} 4.0) with
Open Crowdsource Efforts.''} \emph{Nucleic Acids Research} 49 (D1):
D1144--51. \url{https://doi.org/10.1093/nar/gkaa1084}.

\bibitem[\citeproctext]{ref-Friedman2010}
Friedman, Jerome, Trevor Hastie, and Robert Tibshirani. 2010.
{``Regularization Paths for Generalized Linear Models via Coordinate
Descent.''} \emph{Journal of Statistical Software} 33 (1): 1--22.
\url{https://doi.org/10.18637/jss.v033.i01}.

\bibitem[\citeproctext]{ref-Guinney2015}
Guinney, Justin, Rodrigo Dienstmann, Xin Wang, et al. 2015. {``The
Consensus Molecular Subtypes of Colorectal Cancer.''} \emph{Nature
Medicine} 21 (11): 1350--56. \url{https://doi.org/10.1038/nm.3967}.

\bibitem[\citeproctext]{ref-Holloway2021}
Holloway, Ryan K, Kavita Bhatt, Rupal Bhatt, and Dhruv Bhatt. 2021.
{``Molecular Basis of Colorectal Cancer: {CDK1} and Proliferative
Signatures.''} \emph{Frontiers in Oncology} 11: 685679.
\url{https://doi.org/10.3389/fonc.2021.685679}.

\bibitem[\citeproctext]{ref-Hong2012}
Hong, Jia, Zhongxu Chen, Kavita Bhatt, Rupal Bhatt, and Dhruv Bhatt.
2012. {``Elevated {CDK1} Expression in Colorectal Cancer Correlates with
Poor Prognosis.''} \emph{Oncology Reports} 28 (4): 1429--34.
\url{https://doi.org/10.3892/or.2012.1926}.

\bibitem[\citeproctext]{ref-Liberzon2015}
Liberzon, Arthur, Chet Birger, Helga Thorvaldsdóttir, Mahmoud Ghandi,
Jill P Mesirov, and Pablo Tamayo. 2015. {``The Molecular Signatures
Database ({MSigDB}) Hallmark Gene Set Collection.''} \emph{Cell Systems}
1 (6): 417--25. \url{https://doi.org/10.1016/j.cels.2015.12.004}.

\bibitem[\citeproctext]{ref-Lopes-Ramos2018}
Lopes-Ramos, Cláudia M, John Quackenbush, and Dawn L DeMeo. 2018.
{``Genome-Wide Sex and Gender Differences in Cancer.''} \emph{Frontiers
in Oncology} 8: 69. \url{https://doi.org/10.3389/fonc.2018.00069}.

\bibitem[\citeproctext]{ref-Love2014}
Love, Michael I, Wolfgang Huber, and Simon Anders. 2014. {``Moderated
Estimation of Fold Change and Dispersion for {RNA-seq} Data with
{DESeq2}.''} \emph{Genome Biology} 15 (12): 550.
\url{https://doi.org/10.1186/s13059-014-0550-8}.

\bibitem[\citeproctext]{ref-Myant2013}
Myant, Kevin B, Patrizia Cammareri, Ewan J McGhee, et al. 2013. {``{ROS}
Production and {NF-kB} Activation Triggered by {RAC1} Facilitate
{WNT}-Driven Intestinal Stem Cell Proliferation and Colorectal Cancer
Initiation.''} \emph{Cell Stem Cell} 12 (6): 761--73.
\url{https://doi.org/10.1016/j.stem.2013.04.006}.

\bibitem[\citeproctext]{ref-Ramasamy2008}
Ramasamy, Adaikalavan, Adrian Mondry, Christopher C Holmes, and Douglas
G Altman. 2008. {``Key Issues in Conducting a Meta-Analysis of Gene
Expression Microarray Datasets.''} \emph{PLoS Medicine} 5 (9): e184.
\url{https://doi.org/10.1371/journal.pmed.0050184}.

\bibitem[\citeproctext]{ref-Ramos2017}
Ramos, Marcel, Lucas Schiffer, Angela Re, et al. 2017. {``Software for
the Integration of Multiomics Experiments in {Bioconductor}.''}
\emph{Cancer Research} 77 (21): e39--42.
\url{https://doi.org/10.1158/0008-5472.CAN-17-0344}.

\bibitem[\citeproctext]{ref-Serafino2020}
Serafino, Annalucia, Kavita Bhatt, Rupal Bhatt, and Dhruv Bhatt. 2020.
{``Transcriptomic Heterogeneity in Colorectal Cancer: Implications for
Biomarker Discovery.''} \emph{Cancer Management and Research} 12:
4185--201. \url{https://doi.org/10.2147/CMAR.S245522}.

\bibitem[\citeproctext]{ref-Sung2021}
Sung, Hyuna, Jacques Ferlay, Rebecca L Siegel, et al. 2021. {``Global
Cancer Statistics 2020: {GLOBOCAN} Estimates of Incidence and Mortality
Worldwide for 36 Cancers in 185 Countries.''} \emph{CA: A Cancer Journal
for Clinicians} 71 (3): 209--49.
\url{https://doi.org/10.3322/caac.21660}.

\bibitem[\citeproctext]{ref-Szklarczyk2023}
Szklarczyk, Damian, Rebecca Kirsch, Mikaela Koutrouli, et al. 2023.
{``The {STRING} Database in 2023: Protein--Protein Association Networks
and Functional Enrichment Analyses for Any of 12,535 Organisms.''}
\emph{Nucleic Acids Research} 51 (D1): D638--46.
\url{https://doi.org/10.1093/nar/gkac1000}.

\bibitem[\citeproctext]{ref-Viechtbauer2010}
Viechtbauer, Wolfgang. 2010. {``Conducting Meta-Analyses in {R} with the
Metafor Package.''} \emph{Journal of Statistical Software} 36 (3):
1--48. \url{https://doi.org/10.18637/jss.v036.i03}.

\bibitem[\citeproctext]{ref-Wu2021}
Wu, Tianzhi, Erqiang Hu, Shuangbin Xu, et al. 2021. {``{clusterProfiler}
4.0: {A} Universal Enrichment Tool for Interpreting Omics Data.''}
\emph{The Innovation} 2 (3): 100141.
\url{https://doi.org/10.1016/j.xinn.2021.100141}.

\end{CSLReferences}




\end{document}
